\begin{table}[htbp]
\centering
\ifdefined\cnexperimenttableformat\cnexperimenttableformat\fi
\captionsetup{name=Table}
\caption{CFDepth procedure.}
\label{tab:cfdepth_algorithm}
\small
\setlength{\tabcolsep}{4pt}
\renewcommand{\arraystretch}{1.10}
\begin{tabularx}{\linewidth}{@{}>{\raggedleft\arraybackslash}p{0.035\linewidth}X@{}}
\toprule
\multicolumn{2}{@{}l}{\textbf{Input:} \(\hat{\mathbf M}\): flood map; \(Z\): DEM;} \\
\multicolumn{2}{@{}l}{\hspace{2em}\(d_{\min}\): minimum positive depth; \(D_{\max}\): maximum propagation distance.} \\
\multicolumn{2}{@{}l}{\textbf{Output:} flood-depth estimate \(\hat d\).} \\
\midrule
1: & Align \(\hat{\mathbf M}\) and \(Z\) to a common grid and define flood support \(\Omega_f\). \\
2: & Label the eight-connected components of \(\Omega_f\). \\
3: & Extract the outer boundary and retain candidates that satisfy the slope and local median-deviation criteria; extract the inner boundary by erosion. \\
4: & Pair the outer- and inner-side boundary medians and construct the water-surface anchor \(S_b=\max[(Z_{\mathrm{in}}+Z_{\mathrm{out}})/2,\,Z+d_{\min}]\). \\
5: & For each component, find the nearest reachable anchor along an eight-neighbor path with distance at most \(D_{\max}\). \\
6: & Set \(S^{(0)}=\max[S_b(b(p)),Z+d_{\min}]\) for pixels with a reachable anchor; otherwise use the terrain-floor fallback \(Z+d_{\min}\). \\
7: & Refine the water surface with the \(3\times3\) neighbor mean, relaxation coefficient \(\eta\), fixed anchors, and the terrain floor. \\
8: & Compute \(\hat d=\max[S^{(K)}-Z,0]\) for \(p\in\Omega_f\), and set the remaining pixels to NoData. \\
\bottomrule
\end{tabularx}
\end{table}
